\section{Fibration data and reproducibility}
\label{app:reproducibility}

\subsection{Integral E-string root images}

Across the $C_1$ flop the cubic changes by
$\kappa_{\rm fl}(A,B,C)=\kappa(A,B,C)-(AC_1)(BC_1)(CC_1)$.
Intersecting the flat-model sections with $D_3,D_4$ gives the
following curve pairings in the primitive divisor basis of
Section~\ref{subsec:x9}:
\begin{equation}
\begin{array}{c|rrrrrr}
 &D_2&D_5&D_6&D_7&D_8&E\\\hline
 o_3=-C_1&-1&0&1&1&0&0\\
 o_4=C_2&-1&1&0&1&0&0\\
 s_3&0&0&1&0&1&-1\\
 s_4&0&1&0&0&1&-1\\
 F_3=F_4&1&0&0&0&0&1\\
 r_3=r_4&0&0&0&-1&1&-2
\end{array}
\label{eq:x9-estring-pairing-table}
\end{equation}
The unimodular pairing in Appendix~\ref{app:integral-charges}
makes equality of these rows an integral homology statement.
In particular $s_4-s_3=C_1+C_2$, whereas the common root
$r_i=s_i-o_i-F_i$ is a different class. This distinction is
needed when assigning the elliptic Kaluza--Klein charge.

\subsection{An elliptic K3 in the common slice}

The slice $m_4+2m_1+m_2=0$ of the Newton polytope has $30$
lattice points. In an integral basis its polar has vertices
\begin{equation}
 (-4,-2,-1),\ (0,1,0),\ (1,0,1),\ (1,0,0),\
 (-1,-1,0),\ (0,0,1).
\end{equation}
The two pieces of the subdivision meet in this reflexive
three-dimensional slice. The resulting K3 equation can be
written, in an affine base coordinate $s$, as
\begin{equation}
 a_0w^2+(b_0u^2+b_1uv+b_2v^2)w
 +u(c_0u^3+c_1u^2v+c_2uv^2+c_3v^3)=0.
\label{eq:x9-k3-slice-equation}
\end{equation}
Here $a_0,b_2$ are nonzero constants,
$\deg b_1\leq2$, $\deg b_0\leq4$, and
$c_j\in s\,\CC[s]_{\leq7-2j}$. An exact rational member
is specified by the following coefficient vectors, ordered
by increasing power of $s$:
\begin{equation}
\begin{array}{c|l}
 a_0&(4)\\ b_2&(2)\\ b_1&(3,1,5)\\ b_0&(7,4,4,9,5)\\
 c_3/s&(7,4)\\ c_2/s&(1,8,5,6)\\
 c_1/s&(2,3,1,5,9,4)\\ c_0/s&(5,9,8,6,5,1,3,8)
\end{array}
\label{eq:x9-k3-coefficients}
\end{equation}
Completing the square produces a binary quartic with coefficients
\begin{equation}
\begin{aligned}
 A&=b_0^2-4a_0c_0,& B&=2b_0b_1-4a_0c_1,\\
 C&=b_1^2+2b_0b_2-4a_0c_2,&D&=2b_1b_2-4a_0c_3,&E_4&=b_2^2.
\end{aligned}
\end{equation}
Writing $E_4$ for the quartic coefficient, to distinguish it
from the divisor $E$, its invariants are
\begin{equation}
 I=12AE_4-3BD+C^2,\qquad
 J_4=72ACE_4+9BCD-27AD^2-27E_4B^2-2C^3.
\end{equation}
One normalization of the Jacobian has $f=-I/3$, $g=-J_4/27$.
The polynomial $4I^3-J_4^2$ has degree $24$ and a double
zero at $s=0$. Its quotient by $s^2$ is squarefree and coprime
to $I,J_4$; the fiber at infinity is smooth. Thus this member
has precisely $I_2+22I_1$. This is an exact K3 calculation,
not a construction of the heterotic bundles or of a semistable
total degeneration.

\subsection{Reproducibility and the role of finite checks}

The source distribution contains exact-arithmetic calculations
in \texttt{checks/}, with commands and dependencies documented
in \texttt{checks/README.md}. They reconstruct the fan, Chow
ring, primitive lattices, local smoothing topology, mirror
critical scheme, Laurent mutation, Frobenius coefficients and
fibration data printed here. The explicit global smoothing
family is the one constructed in \cite{Wuebben:2026smoothing}.

The period statements do not follow from a finite expansion.
The coefficient identities, analytic continuation estimates
and integral normalization used in their proofs appear in
Section~\ref{sec:quantum} and Appendix~\ref{app:mirror-periods}.
The hypermultiplet threshold check differentiates the singular
prepotential, verifies the four unit symplectic shifts and their
combined Smith invariants, and checks that all neutral charge
projections vanish. It also verifies the algebraic reconstruction
of the supergravity period matrix. These are consistency checks
of the printed charge data, not independent proofs of its geometric marking.
Exact low-degree series comparisons test the period formulas; they
are not substituted for their all-orders arguments. Numerical
central-charge evaluations and the separate magnetic-sheaf
calculations are discussed in the companion note
\cite{Wuebben:2026magnetic}.
